% Written by Claude Fable 5.2 (AI, orchestrator), 2026-11-14. A second proof ("every node is worth -p").
% Lemmas checked numerically in experiments/E004_accounting_proof/claude_check.py. Not human-verified.
% A separate draft of the same idea by a Kimi worker is in alt_proof_accounting.tex; this note supersedes it.
\documentclass[11pt]{article}
\usepackage[margin=1in]{geometry}
\usepackage{amsmath,amssymb,amsthm}
\usepackage{enumitem}
\usepackage[colorlinks=true,urlcolor=blue,linkcolor=blue]{hyperref}
\newtheorem{theorem}{Theorem}
\newtheorem{lemma}{Lemma}
\newtheorem{corollary}{Corollary}
\newtheorem{proposition}{Proposition}
\theoremstyle{remark}
\newtheorem{remark}{Remark}
\newcommand{\Z}{\mathbb Z}
\newcommand{\F}{\mathfrak F}
\newcommand{\C}{C_n^{(q)}}
\setlength{\parskip}{0.4em}
\setlength{\parindent}{0pt}

\title{Every node is worth $-p$:\\ a second proof by counting labelled trees}
\author{Mighty Math Magicians, Caltech Mathathon 2026}
\date{14 November 2026\\\small AI-written draft; lemmas checked by computer; not yet verified by a human}

\begin{document}
\maketitle

\noindent\fbox{\parbox{\dimexpr\linewidth-2\fboxsep-2\fboxrule}{%
\textbf{Status.} This note gives a second, shorter route to the theorem of the main paper
(\texttt{paper/main.pdf}). It is not independent of that paper: it relies on the results
listed in Section~1, in particular on the carry proposition. And it \emph{does} use the
valuation of $m!$: Lemmas~\ref{L3} and~\ref{L4} together are Legendre's formula, proved here by
the same kind of symmetry argument rather than quoted. Whether that satisfies ``without
relying on factorial valuations'' is for the reader to judge. The main paper never mentions
$m!$.}}

\section{Statement and what is borrowed}

Let $p$ be prime, $k\ge1$, $q=p^k$, and let $\C$ be the number of $q$-ary trees with $n$
internal nodes, equivalently of $p$-ary trees with $m=(q-1)n+1$ leaves all of whose leaves have
depth divisible by $k$ (``$k$-blocked''). Write $m=\sum_j m_jp^j$ in base $p$ and
$s=\sum_j m_j$. Put
\[
 N(m)=\frac{m-1}{p-1},\qquad \nu(m)=\frac{m-s}{p-1},\qquad t=N(m)-\nu(m)=\frac{s-1}{p-1}.
\]
Here $N(m)$ is the number of internal nodes of any $p$-ary tree with $m$ leaves, and $\nu(m)$
is the number of internal nodes of the \emph{digit forest} of $m$: $m_j$ complete $p$-ary
trees of height $j$ for each $j$ (a complete tree of height $j$ has $(p^j-1)/(p-1)$ internal
nodes, and $\sum_j m_j\frac{p^j-1}{p-1}=\frac{m-s}{p-1}$).

\begin{theorem}\label{thm}
$\displaystyle \C\cdot\prod_j m_j!\;\equiv\;(-p)^{t}\pmod{p^{t+1}}.$
In particular $v_p(\C)=t$ and $\C/p^t\equiv(-1)^t/\prod_jm_j!\pmod p$.
\end{theorem}

We use the following results of the main paper.
\begin{enumerate}[label=(\Alph*)]
\item \emph{Blocked form and class size.} Lemma ``Blocked form'': $q$-ary trees with $n$
      internal nodes correspond to $k$-blocked $p$-ary trees with $m$ leaves. Lemma ``Class
      size'', in its version for trees with labelled leaves given by Lemma ``Plane
      representatives'': a rotation class of plane labelled $p$-ary trees has $p^{d}$ elements,
      $d$ being the number of asymmetric nodes; rotation-equivalent trees have the same leaf
      depths and the same labels; and the classes of plane trees with distinctly labelled
      leaves are the labelled cyclic trees. Lemma ``Skeleton nodes are asymmetric'': with
      pairwise distinct labels every internal node is asymmetric.
\item \emph{The carry.} Proposition ``A carry collapses a cherry'': if $T(\mathbf a)$ denotes
      the number of cyclic $p$-ary trees with distinctly labelled leaves, $a_c$ of them of
      class $c\in\Z/k$, a leaf of class $c$ lying at depth $\equiv-c\pmod k$, then
      $T(\mathbf a)\equiv-T(\mathbf a-p\,e_c+e_{c+1})\pmod p$ whenever $a_c\ge p$. Lemma
      ``Terminal vectors'': if every $a_c<p$ then $T(\mathbf a)=0$ unless $\mathbf a=e_0$, and
      $T(e_0)=1$. Lemma ``Automorphisms of cyclic trees''\,(i): an automorphism fixing every
      leaf is the identity.
\end{enumerate}
We also use the elementary fact that if $\Z_p$ acts on a finite set $X$ then
$|X|\equiv|\mathrm{Fix}(X)|\pmod p$, because every orbit has $1$ or $p$ elements, $p$ being
prime.

The next lemma is the proposition ``Iterated carries'' of the main paper, stated for an
arbitrary vector; its proof there is written for the vector of digits of $m$, and we repeat it.

\begin{lemma}[Iterated carries, general form]\label{L0}
Let $\mathbf a=(a_0,\dots,a_{k-1})$ be non-negative integers with $\sum_ca_c\ge1$, and if
$q>2$ assume $\sum_ca_cp^c\equiv1\pmod{q-1}$. Then $(\sum_ca_c-1)/(p-1)$ is an integer and
\[
 T(\mathbf a)\equiv(-1)^{(\sum_ca_c-1)/(p-1)}\pmod p .
\]
\end{lemma}
\begin{proof}
As long as some $a_c\ge p$, replace $\mathbf a$ by $\mathbf a-p\,e_c+e_{c+1}$ (a \emph{carry});
by (B) this changes $T$ by the factor $-1$ modulo $p$. A carry lowers $\sum_ca_c$ by $p-1$ and
needs $p$ labels, so the process stops and $\sum_ca_c\ge1$ throughout. Let $\mathbf b$ be the
vector at the end, with all $b_c<p$.
If $q>2$: the quantity $W=\sum_ca_cp^c$ is unchanged modulo $q-1$ by a carry (it changes by $0$
if $c<k-1$ and by $-p^k+1$ if $c=k-1$), so $W(\mathbf b)\equiv1$; also
$0<W(\mathbf b)\le(p-1)\sum_{c<k}p^c=q-1$; hence $W(\mathbf b)=1$ and $\mathbf b=e_0$.
If $q=2$: there is one class and a carry lowers $a_0$ by one, so the process ends at $a_0=1$.
In both cases $\mathbf b=e_0$, $T(e_0)=1$, and the number of carries is
$(\sum_ca_c-1)/(p-1)$.
\end{proof}

\section{Trees: each of the $N(m)$ nodes is worth $-p$}

\begin{lemma}\label{L1}
Let $L(m)$ be the number of cyclic $k$-blocked $p$-ary trees whose $m$ leaves are labelled
$1,\dots,m$. Then
\[
 m!\cdot\C=p^{N(m)}\cdot L(m),\qquad L(m)\equiv(-1)^{N(m)}\pmod p .
\]
\end{lemma}
\begin{proof}
By (A) there are $m!\cdot\C$ plane $k$-blocked trees with leaves labelled $1,\dots,m$; they
fall into rotation classes, each consisting of $k$-blocked trees with the same labels and each
a labelled cyclic tree; and each class has $p^{d}$ elements where $d$ is the number of
asymmetric nodes. With distinct labels all $N(m)$ internal nodes are asymmetric. This gives
the first identity.

All leaves are at depth $\equiv0\pmod k$, so in the notation of (B) they have class $0$ and
$L(m)=T(m\,e_0)$. Since $m\ge1$ and $m\equiv1\pmod{q-1}$, Lemma~\ref{L0} applies with
$\mathbf a=m\,e_0$ and gives $L(m)\equiv(-1)^{(m-1)/(p-1)}$.
\end{proof}

\section{Labellings: each of the $\nu(m)$ nodes of the digit forest is worth $-p$}

In this section $m\ge0$ is an arbitrary integer (no congruence condition), with base-$p$
digits $m_j$ and digit sum $s$; $\nu(m)=(m-s)/(p-1)$ as before.

A \emph{digit-forest structure} on $\{1,\dots,m\}$ is a partition of this set into $m_j$
blocks of size $p^j$ for each $j$, each block arranged as the set of leaves of a cyclic
complete $p$-ary tree of height $j$ (for $j=0$: a single point). Let $\F(m)$ be the set of
these structures; $\F(0)$ has one element.

\begin{lemma}\label{L3}
$m!=p^{\nu(m)}\cdot\prod_jm_j!\cdot|\F(m)|$.
\end{lemma}
\begin{proof}
The symmetric group $S_m$ acts on $\F(m)$ by relabelling points, transitively: two structures
have the same number of trees of each height, so some bijection of points carries one to the
other. A permutation fixes a structure exactly when it maps each tree to a tree of the same
height by an isomorphism of cyclic trees. The trees of height $j$ can be permuted among
themselves in $m_j!$ ways, and a cyclic complete tree of height $j$ has
$p^{(p^j-1)/(p-1)}$ automorphisms: an automorphism is a free choice of a rotation at each of
its $(p^j-1)/(p-1)$ internal nodes (by induction on $j$: a rotation at the root, then an
automorphism of each of the $p$ subtrees), and distinct automorphisms induce distinct
permutations of the leaves, because an automorphism fixing every leaf is the identity (B).
So the stabilizer has order
$\prod_j m_j!\,p^{m_j(p^j-1)/(p-1)}=p^{\nu(m)}\prod_jm_j!$, and the orbit--stabilizer count
gives the identity.
\end{proof}

\begin{lemma}\label{L4}
$|\F(m)|\equiv(-1)^{\nu(m)}\pmod p$.
\end{lemma}
\begin{proof}
Write $m=pa+m_0$ with $a=\lfloor m/p\rfloor$. In a digit-forest structure the $m_0$ points
that are trees of height $0$ form a subset $R$, and the other $pa$ points are grouped into $a$
\emph{cherries}: the sets of $p$ leaves below the lowest internal nodes, each with its cyclic
order. Contracting every cherry to a point leaves a digit-forest structure on the set of the
$a$ cherries (the digits of $a$ are $m_1,m_2,\dots$, and a tree of height $j\ge1$ becomes one
of height $j-1$). Conversely $R$, the cherries with their cyclic orders, and a digit-forest
structure on the set of cherries determine the structure. A way of grouping $pa$ points into
$a$ cyclically ordered $p$-sets is the same thing as a permutation of those points all of
whose cycles have length $p$; let $K(a)$ be the number of such permutations. Hence
\begin{equation}\label{rec}
 |\F(m)|=\binom{m}{m_0}\cdot K(a)\cdot|\F(a)| .
\end{equation}

\emph{$\binom{pa+m_0}{m_0}\equiv1\pmod p$ for $0\le m_0<p$.} Let $\Z_p$ act on the
$m_0$-subsets of $\{1,\dots,pa+m_0\}$ through the permutation $(1\,2\,\cdots\,p)$. A fixed
subset contains all or none of $1,\dots,p$; having fewer than $p$ elements it contains none.
So $\binom{pa+m_0}{m_0}\equiv\binom{p(a-1)+m_0}{m_0}$, and after $a$ steps we reach
$\binom{m_0}{m_0}=1$.

\emph{$K(a)\equiv(-1)^a\pmod p$.} Let $\Z_p$ act on the permutations of $\{1,\dots,pa\}$ with
all cycles of length $p$ by conjugation with $c=(1\,2\,\cdots\,p)$. A fixed permutation $\pi$
commutes with $c$. Then $\pi$ maps the set $\{1,\dots,p\}$ of points moved by $c$ to itself,
and if $\pi(1)=c^u(1)$ then $\pi(c^i(1))=c^i(\pi(1))=c^{i+u}(1)$ for all $i$, so $\pi$ acts on
$\{1,\dots,p\}$ as $c^u$; $u\ne0$ because $\pi$ has no fixed point, and each $c^u$ with
$u\ne0$ is a single $p$-cycle because $p$ is prime. The rest of $\pi$ is any such permutation of the other
$p(a-1)$ points. So $K(a)\equiv(p-1)K(a-1)\equiv-K(a-1)$, and $K(0)=1$.

By \eqref{rec} and induction, $|\F(m)|\equiv(-1)^a|\F(a)|\equiv(-1)^{a+\nu(a)}$, and
$a+\nu(a)=\nu(m)$: indeed $pa=m-m_0$ and the digit sum of $a$ is $s-m_0$, so
$a+\frac{a-(s-m_0)}{p-1}=\frac{pa+m_0-s}{p-1}=\frac{m-s}{p-1}$. (The $a$ cherries are the
internal nodes of the digit forest at the lowest level.)
\end{proof}

\begin{corollary}[Legendre's formula, with its residue]
For every integer $m\ge0$, $v_p(m!)=\nu(m)$ and
$m!/p^{\nu(m)}\equiv(-1)^{\nu(m)}\prod_jm_j!\pmod p$.
\end{corollary}
\begin{proof}
Immediate from Lemmas~\ref{L3} and~\ref{L4}, since $\prod_jm_j!$ is prime to $p$.
\end{proof}

\section{Proof of the theorem}

Divide the identity of Lemma~\ref{L1} by that of Lemma~\ref{L3}:
\[
 \C\cdot\prod_jm_j!\cdot|\F(m)|=p^{\,N(m)-\nu(m)}\cdot L(m)=p^{\,t}\cdot L(m).
\]
By Lemma~\ref{L4}, $|\F(m)|$ is prime to $p$, so $p^t$ divides $\C\prod_jm_j!$; write
$\C\prod_jm_j!=p^tx$. Then $x\,|\F(m)|=L(m)$, and modulo $p$ this reads
$x\,(-1)^{\nu(m)}\equiv(-1)^{N(m)}$, that is, $x\equiv(-1)^t$. Hence
$\C\prod_jm_j!\equiv(-1)^tp^t\pmod{p^{t+1}}$. Since $\prod_jm_j!$ is prime to $p$ (every factor
is less than $p$), $v_p(\C)=t$ and the residue follows. \qed

\section{What it means}

\textbf{The accounting.} A tree with labelled leaves has $N(m)$ internal nodes, and each
contributes a factor $p$ (the $p$ rotations at that node give different labelled trees) and a
sign $-1$ (from the carry: $p-1\equiv-1$). The $m!$ labellings contain $\nu(m)$ such factors
of their own, one for each internal node of the digit forest, and again each comes with a
sign $-1$. The difference, $t$ nodes, is what is left in the number of unlabelled trees.

\textbf{Relation to the main proof.} There, a minimal tree is a skeleton of $t$ asymmetric
nodes carrying atoms, and the atoms are the digit forest of $m$. So the $N(m)$ nodes of a
minimal tree split as $\nu(m)$ nodes inside atoms plus $t$ skeleton nodes: the same
subtraction, seen inside one tree.

\textbf{Sand.} Consider grains on sites with the rule: \emph{$p$ grains on a site topple into
one grain on the next site.}
\begin{proposition}
\textup{(a)} On the half-line of sites $0,1,2,\dots$, starting from $m$ grains on site $0$, every
sequence of topplings ends, after exactly $\nu(m)$ topplings, in the configuration given by
the base-$p$ digits of $m$.
\textup{(b)} On a ring of $k$ sites, starting from $m\ge1$ grains on site $0$ with
$m\equiv1\pmod{p^k-1}$, every sequence of topplings ends, after exactly $N(m)$ topplings, in
a single grain on site $0$.
Consequently $v_p(\C)=t$ is the number of topplings on the ring minus the number on the line.
\end{proposition}
\begin{proof}
Each toppling removes $p-1$ grains, so every sequence ends, and it ends when every site holds
fewer than $p$ grains. (a) The quantity $\sum_jg_jp^j$ ($g_j$ grains on site $j$) is unchanged
by a toppling and equals $m$ at the start, so the final configuration is the base-$p$
expansion of $m$, with $s$ grains; the number of topplings is $(m-s)/(p-1)$. (b) The quantity
$\sum_cg_cp^c$ modulo $p^k-1$ is unchanged and is $\equiv1$; a final configuration has
$0<\sum_cg_cp^c\le p^k-1$, hence $\sum_cg_cp^c=1$, a single grain on site $0$ (for $q=2$ there
is one site and the count of grains simply drops to $1$); the number of topplings is
$(m-1)/(p-1)$.
\end{proof}
In the proofs above a toppling on the ring is a cherry collapsing in a labelled tree (Lemma
\ref{L1}), and a toppling on the line is a cherry of the digit forest being contracted (Lemma
\ref{L4}); each contributes the factor $-p$. For $q>2$, if $p^k-1$ divides $m\ge1$ the
ring ends instead in the full configuration $(p-1,\dots,p-1)$; this case does not occur for
$m=(q-1)n+1$. (For $q=2$ every $m\ge1$ ends in a single grain.)

\end{document}
